\paragraph{A finite divisor on the vertex quotient.}
The following equation on the genus-seven quotient gives a finite
index set for testing the remaining common-line directions. On the
rank-five bundle \(W\) above, take the two vertex-invariant linear forms
\[
 H(X)=X_{02}+X_{12}+X_{23},\qquad D(X)=X_{00}+X_{11}+X_{33}.
\]
Both kill \(B\). Indeed \(H(B)=m_2(B)=0\), and if \(T\) is the sum of its
six edges, the corrected diagonal entries are \(T/2-2m_i(B)\).
The sum of the three outer diagonal entries is \(3T/2-4T=0\) at five.
Thus the two forms are defined on \(W_B\).

Write \(b_j\) for the symmetric bilinear form with
\(b_j(x,x)=Q_j(x)\). Wedge the four covectors
\(b_2(x,-), b_3(x,-), b_4(x,-), D\) to get a cofactor vector
\(c(x)\in W\otimes\det W^*\otimes\mathcal O_C(3)\), cubic in \(x\).
Set \(S(x)=Q_2(c(x))\) and form the homogeneous resultant
\[
         R=\operatorname{Res}_x(H,Q_2,Q_3,Q_4,S).
\]
If a two-space \(P\subset W_B\) is common isotropic, take nonzero
\(x\in P\cap\ker H\) over the algebraic closure. The three polar rows
annihilate \(P\). At rank three their common kernel is \(P\), hence
contains the cofactor. At smaller rank the cofactor is zero. In either
case \(S(x)=0\), so \(R(B)=0\). At \(B=(1,1,3,0,3,4)\) the checker uses
\(D,H\) as the final two coordinates of \(W_B\). The quartic reductions
exclude points at infinity in \(H=0\), and the reduced sextic has a
verified inverse on its affine chart. Hence \(R(B)\ne0\).

The five degrees are \((1,2,2,2,6)\), with product 48. The resultant
has degree 48 divided by the corresponding degree in each coefficient
block \cite{dandreaDickensteinResultant}. An invertible change of
variables preserves the universal resultant hypersurface, so it acts on
the resultant by a determinant character. Substitution of a scalar
matrix in the five variables gives exponent 48 on its determinant.
Writing \(M=\det W^*\), the sextic
has coefficient twist \(M^2\otimes\mathcal O_C(6)\). Altogether the
resultant takes values in \(M^{64}\otimes\mathcal O_C(120)\).
Equivariantly \(M=\det(V^*)\otimes\mathcal O_C(5)\), where \(V\) is the
nine-dimensional tangent space before imposing the three polars and
quotienting by \(B\). Its determinant character has square one.
Thus \(R\) is an invariant section of \(\mathcal O_C(440)\).
As \(Y=\sum_i y_i\) is also invariant and nonzero at this test point,
\(h=R/Y^{440}\) is a nonzero rational function on the smooth quotient.
Its pullback has at most \(24\cdot440\) zeros counted with multiplicity.
The degree-six quotient map gives \(\deg (h)_0\le 1760\).
Every common-line direction on the face has \(Y\ne0\).

\begin{theorem}[Decomposition of the common-line incidence on the fixed face]
\label{thm:common-line-decomposition}
Work in characteristic five on the face \([l:m]=[0:1]\).
Let \(U\) be the invariant pure chart with distinct incident edges,
\(w=1\), all six edges nonzero, and \(u(v-3)AY\ne0\), where \(A\)
is the leading coefficient in the first subresultant. It is an open
of the genus-seven quotient with coordinates \(p,q,u,v\). Every
geometric common-line direction on the face lies above \(U\).
Let \(T\to U\) order the three roots of \(X^3+pX-q\).
There is an explicit scheme \(\mathcal F\to U\) of lines in the
descended net. The incidence scheme \(\mathcal I\) of labeled
directions with common lines is the pullback
\[
                \mathcal I\simeq T\times_U\mathcal F .
\]
The projection of \(\mathcal F\) has geometric support in
\(U\cap(h)_0\). At each closed point \(s\) of this candidate divisor
let \(K=\kappa(s)\), let \(e\) be the splitting degree of the incident
cubic over \(K\), and let \(K_e/K\) be the degree-\(e\) extension.
There is an isomorphism of \(K\)-schemes
\[
     \mathcal I_s\simeq
        \coprod_{j=1}^{6/e} (\mathcal F_s\times_K K_e).
\]
In particular \(\mathcal I_s\) is geometrically empty exactly when
\(\mathcal F_s\) is. The fiber \(\mathcal F_s\) retains its scheme structure.
\end{theorem}

\begin{proof}
The exclusions above remove the complement of \(U\) from the possible
incidence. Over its coordinate algebra the ordered-root algebra has
presentation
\[
  x^3+px-q=0,\qquad y^2+xy+x^2+p=0,\qquad z=-x-y .
\]
It is free of rank six by the two monic equations. Inverting the
discriminant makes it an étale \(S_3\)-torsor.
For a tangent vector, evaluate three polynomials of degree at most two
at the ordered roots, one polynomial each for the incident edges, the
opposite edges, and the three outer diagonal entries. The Vandermonde
matrix is invertible. The tangent sum determines the central diagonal.
Adding a multiple of \(B\) shifts the linear coefficient of the
incident polynomial by that multiple. Setting the coefficient to zero
gives eight invariant coordinates. The three polar equations in these
coordinates have constant rank three, since they do after the faithfully
flat extension ordering the roots. Denote their rank-five kernel by
\(E\). In \(\operatorname{Gr}_U(2,E)\) let \(\mathcal F\) be the zero scheme
of the restrictions of the three descended quadrics to the universal
two-space. In a frame these are nine Gram pairings. After pullback to
\(T\) the interpolation is an invertible change of coordinates and
these polynomials are precisely the three original restricted quadrics.
The construction commutes with base change, giving the claimed scheme
isomorphism. The resultant gives the support bound.

Over the residue field, Frobenius permutes the three distinct roots.
On the set of their six orderings it acts freely through the cyclic
group it generates in \(S_3\). All its orbits have size \(e\).
The equivalence between finite étale algebras and finite Galois sets
\cite{stacksEtalePoint} now gives
\[
\begin{array}{c|c|c}
 \text{factor degrees of the cubic over }K & e & \mathcal O(T_s)\\ \hline
 1,1,1 &1&K^6\\
 1,2   &2&K_2^3\\
 3     &3&K_3^2
\end{array}
\]
Pullback gives the formula even for nonreduced or positive-dimensional
line fibers. To retain infinitesimal variation of the direction itself,
let \(Z\subset U\) be the scheme-theoretic image of \(\mathcal F\to U\).
Its support is finite, hence it is a disjoint union of local Artin
schemes. Over each such local ring the same \(6/e\) finite étale
factors lift uniquely up to isomorphism from the residue field
\cite{stacksHenselianEtale}. Pulling back \(\mathcal F\) along those
factors gives the decomposition over the thickened image as well.
\end{proof}

For the degree-\(n\) extension \(K_n/K\), the same formula gives
\(\#\mathcal I_s(K_n)=6\#\mathcal F_s(K_n)\) if \(e\mid n\), and
zero otherwise. Over the algebraic closure every line has six labelings.
The degree bound also gives a finite search independent of an expansion
of \(h\): enumerate the closed points of \(U\) of degree at most 1,760
over \(\mathbb F_5\), recovering \(p,q\) from \(u,v\). Larger-degree points
cannot belong to the support. At each candidate \(s\), the
ten row-echelon charts decide geometric existence in \(\mathcal F_s\).
On a proper chart the field search recovers a line. Over a common
field of definition, evaluating its two rows at the ordered roots
recovers a frame \((B,X_1,X_2)\) and permits direct verification in
the original permanent. If a full constant-permanent plane is required,
impose the remaining cubic and quartic tangent coefficients on the
whole line family. Vanishing of all nonconstant coefficients is
preserved under a change of frame, adding multiples of \(B\) to its
representatives, and permuting the vertices. It therefore defines a
closed subfamily with the same ordering decomposition. For an individual
recovered line these extra coefficients test that particular lift.
The factorization of \((h)_0\) and the list of its nonempty line
fibers remain undetermined.

\paragraph{A complete test over quadratic slopes.}
For each of the ten monic irreducible quadratics \(v^2+sv+t\) over
\(\mathbb F_5\), reducing the invariant plane equation gives
\(a(u)+b(u)v\). In every case \(a,b\) are coprime and
\(a^2-sab+tb^2\) is squarefree of degree 24. Thus \(u\) separates
the geometric points and \(v=-a/b\) recovers the slope. The lifts
satisfy the original pure equations with separable incident cubic.
Across the ten quadratics there are 30 closed quotient points.
Three degree-two points lie on the excluded pencils. At each of the
other 27 the degree-nine Jacobian rank is 54, above the common-line
ceiling of 52. Thus every line fiber over these slopes is geometrically
empty. The ordering algebra has splitting degree one, two and three
at seven, sixteen and seven closed points, respectively. Those give
104 closed labeled directions, comprising 1,440 geometric directions
before testing for a line. With the prime-field slope calculation,
this excludes every \(v\in\mathbb F_{25}\) throughout the eligible chart.
